% Proof-level conductor-scope repair for the next arithmetic gate.
% Run JAC1964-20261121; /root/h9_height; GPT-6 Astra Pro.
% Coordinate with AR19 (Katz value); this is H9-T15, not a Chow-rank theorem.

\subsection*{H9-T15. The conductor-general rank-one implication for the Fermat character}

\paragraph{Why a separate deduction is needed.}
Castella's \emph{Tamagawa number conjecture for CM modular forms and
Rankin--Selberg convolutions}, manuscript of 14 February 2025,
introduces in Section 5, printed p.~21, the standing hypothesis
$\mathfrak d_K\Vert\mathfrak c$. The displayed Theorem 5.2.1,
printed p.~23, does not repeat that hypothesis. It is therefore
inappropriate to disregard it merely by reading the theorem display.
For the Fermat character its conductor is $(1+i)^4$, while the
different of $\mathbf Q(i)$ is $(1+i)^2$; exact divisibility fails.
The argument below instead uses the arbitrary-conductor reciprocity
theorem of Section 4 and the Euler-system bound. It records exactly
the portion of the proof of Theorem 5.2.1 that does not use the
Section 5 good-pair or conductor hypothesis.

\begin{proposition}[H9-T15]
Let $K_0=\mathbf Q(i)$, let $p=5$, and let
$\chi=\psi^3\mu=(\psi')^3$ be the ideal-valued Hecke character of
the Fermat rank-two motive, in the notation of AR-11 and BDP
Section 5.2. Write $\lambda$ for this same character in Castella's
infinity-type convention, so that it has type $(-3,0)$ and central
character $\eta_{K_0}$. Put
\[
 \varphi=\lambda\mathbf N^{-2},\qquad
 V=V_\lambda(2).
\]
Thus $\varphi$ has type $(-1,2)$ and is self-dual in Castella's
Definition 2.2.2. Fix one embedding $K_0\hookrightarrow\mathbf Q_5$
and use the associated prime $v$ and its conjugate $\bar v$ as in
that source. Let $\Phi/\mathbf Q_5$ contain the character values
and let $T\subset V$ be a stable rank-one lattice over its integers
$\mathcal O$. One can take $\Phi=\mathbf Q_5$ in this example.

If
\[
                  \mathcal L_v(\varphi)(0)\ne0,              \tag{H9.56}
\]
then
\[
 \dim_\Phi\operatorname{Sel}_{\rm BK}(K_0,V)=1.             \tag{H9.57}
\]
This implication requires no exact-divisibility hypothesis at $2$.
It is a statement about this specified Selmer group; it does not
identify its dimension with a rational Chow rank.
\end{proposition}

\begin{proof}
We use the conventions of Castella throughout. In particular the
cyclotomic character has Hodge--Tate weight $+1$, and the two
Hodge--Tate weights of $V$ are $-1$ at $v$ and $2$ at $\bar v$.
Self-duality identifies its Tate dual $V^*(1)$ with the conjugate
rank-one representation $V^\tau$. Neither representation is
trivial globally or at either prime over $5$. Their local
Hodge--Tate weights are nonzero, which also proves
\[
 H^0(K_0,V)=H^0(K_0,V^*(1))=0,
 \qquad H^0((K_0)_w,V^*(1))=0\quad(w\mid5).                \tag{H9.58}
\]

\emph{The nonzero class and its local conditions.}
Castella's Theorem 4.2.1, printed p.~16, is stated for an arbitrary
self-dual Hecke character. Its elliptic-unit construction starts
with any prime-to-$p$ conductor ideal; it has no requirement that
the conductor have exact valuation equal to the different at a
ramified prime. Apply this construction to $\varphi$ and its
conjugate. Let
\[
 c\in H^1_{\rm Iw}(K_\infty,T),\qquad
 c_0\in H^1(K_0,T)
\]
be the normalized twisted elliptic-unit class and its trivial
anticyclotomic specialization. The specialization of the Coleman
maps, which interpolate the Bloch--Kato logarithm and dual
exponential in Theorem 4.2.1, gives the two consequences used in
the first half of Castella's printed p.~24 proof:
\[
 \mathcal L_v(\varphi)(0)\ne0
       \ \Longrightarrow\ \operatorname{loc}_{\bar v}(c_0)\ne0,
                                                               \tag{H9.59}
\]
and
\[
 \mathcal L_v(\varphi^\tau)(0)=0
       \ \Longrightarrow\ \operatorname{loc}_{v}(c_0)=0.
                                                               \tag{H9.60}
\]
These are the reciprocity and specialization statements of the
Section 4 theorem, not an invocation of the Section 5 conductor
hypothesis. Conjugating the two primes gives
$\mathcal L_{\bar v}(\varphi^\tau)(0)=\mathcal L_v(\varphi)(0)$
up to the compatible nonzero period unit, explaining the prime
$\bar v$ in (H9.59). The local weights $-1,2$ avoid the exceptional
weight-zero and weight-one cases; after inverting $5$ the relevant
logarithm or dual exponential identifies the one-dimensional local
space in question with its de Rham target.

The conjugate character has type $(2,-1)$, in the interpolation
range of Castella Theorem 2.2.1, printed p.~7. The central root
number of $\lambda$ is $-1$ by the Fermat character computation
in AR-11, so $L(\lambda,2)=L(\lambda^\tau,2)=0$.
Interpolation therefore gives
$\mathcal L_v(\varphi^\tau)(0)=0$, and (H9.60) applies.
Twisted elliptic units are unramified in the required rational
local condition away from $5$. At $5$, Lemma 3.4.2, printed p.~12,
identifies the Bloch--Kato local space with $0$ at $v$ and with
all of $H^1((K_0)_{\bar v},V)$ at $\bar v$. Consequently
\[
       0\ne c_0\otimes1
             \in\operatorname{Sel}_{\rm BK}(K_0,V).        \tag{H9.61}
\]

\emph{The Euler-system bound with the ramified prime retained.}
For clarity, the Euler system is taken over the ray-class tower
containing $K_0(\mathfrak c5^\infty)$ and all auxiliary ray-class
extensions at primes away from $6\mathfrak c5$, where
$\mathfrak c=(1+i)^4$. One starts with the elliptic units for
$\mathcal O(1)$ and twists by the rank-one character
$T\otimes\mathcal O(-1)$. This character factors through a finite
extension of the compositum of the $\mathbf Z_5$-extensions of
$K_0$, inside the displayed ray-class tower.

Rubin's twisting theorem constructs the full twisted Euler system
whose normalized base specialization is $c_0$: in the directly read author draft
\emph{Euler systems}, Chapter VI, Definition 3.1 and Theorem 3.5,
printed pp.~92--93, the new non-$p$ conductor is simply adjoined
to the excluded-prime ideal. Thus the prime above $2$ is among
the excluded bad primes. No constraint on its conductor exponent
is used. This is the same twisting input cited in Castella's
printed p.~24 proof as the published Rubin Theorem 6.4.1.
The harmless normalization of the elliptic units is the invertible
auxiliary-ideal factor in Castella Theorem 4.2.1: its proof explicitly
chooses $\mathfrak a$ prime to $6\mathfrak cp$ such that
$\sigma_{\mathfrak a}-\varphi^\tau(\sigma_{\mathfrak a})$
is invertible in the anticyclotomic Iwasawa algebra. More precisely,
apply Rubin to the \emph{unnormalized} full twisted elliptic-unit
Euler system. Its base class is the nonzero unit multiple of $c_0$
obtained by specializing that auxiliary-ideal factor. Thus one need
not invert the anticyclotomic normalization simultaneously on all
auxiliary ray extensions; its base specialization already supplies
the nonzero Euler-system class required below.

Rubin's Euler-system bound now gives
\[
           \operatorname{Sel}_{\rm str}(K_0,V^*(1))=0.
                                                               \tag{H9.62}
\]
The exact primary statement is Chapter II, Theorem 2.3, printed
p.~24 of that draft (published Theorem 2.2.3): a nonzero
Euler-system class for a nontrivial representation satisfying
$\operatorname{Hyp}(K_0,V)$ makes the strict dual torsion Selmer
group finite, and hence its rational space zero. Its hypotheses
hold for this rank-one representation: choose $\tau=1$ in the
rank-one-quotient condition, and irreducibility is automatic.
The same page's Remark 2.4 explicitly records this rank-one case.
The nontriviality condition follows also from (H9.58).

\emph{The dimension calculation.}
Let $\operatorname{Sel}_{\rm rel}(K_0,V)$ impose the Bloch--Kato
unramified condition away from $5$ and impose no condition at the
two primes over $5$. Poitou--Tate duality and the global Euler
characteristic give, using (H9.58),
\[
 \dim_\Phi\operatorname{Sel}_{\rm rel}(K_0,V)
    =d^-(V)+
       \dim_\Phi\operatorname{Sel}_{\rm str}(K_0,V^*(1)).    \tag{H9.63}
\]
Here $d^-(V)=1$, because $K_0$ has one complex place and $V$ has
rank one. A complete displayed derivation of (H9.63) is
Loeffler--Zerbes, \emph{Euler systems with local conditions},
arXiv:1710.04956v2, Proposition 2.8 and its proof, printed p.~3.
The only suppressed terms in (H9.63) are the global and local
$H^0$ terms explicitly checked in (H9.58).

By (H9.62), the relaxed space in (H9.63) has dimension one.
Its Bloch--Kato subspace contains the nonzero vector (H9.61).
It is therefore one-dimensional, which proves (H9.57).
\end{proof}

\paragraph{Scope of this check.}
This is a proof-level application of the arbitrary-conductor
reciprocity construction and rank-one Euler-system machinery.
It repairs the source-scope issue caused by the standing
Section 5 assumption; it does not claim that assumption was
explicitly withdrawn in the source. The hypothesis (H9.56)
remains a separate Katz-value calculation, assigned to AR19.
Neither a nonzero complex derivative alone nor the BDP
Rankin--Katz formula outside its stated ramification hypotheses
is used to prove it here. No scalar-extended Abel--Jacobi
injectivity, rational Chow-generation theorem, or positivity on
the full arithmetic complementary Chow group follows just from
(H9.57).

\paragraph{Primary sources directly inspected.}
Castella: \url{https://web.math.ucsb.edu/~castella/Katz.pdf},
Section 4.2 and printed p.~24 (the latter also visually inspected).
Rubin author draft:
\url{https://swc-math.github.io/aws/1999/99RubinES.pdf},
Chapter II pp.~23--24 and Chapter VI pp.~92--94.
Loeffler--Zerbes:
\url{https://discovery.ucl.ac.uk/id/eprint/10067153/1/1710.04956v2.pdf},
Theorem 2.5, Proposition 2.8 and its proof, printed pp.~2--3.
The general Euler-system theorems are used as established inputs;
they have not been reproved or formally verified in this run.
